\section{Introduction}

Neural computation is increasingly deployed to hardware under tight constraints: latency budgets, deterministic deadlines, and limited numeric precision. Yet the dominant way these systems are validated remains fundamentally informal: a software model is treated as ``the truth,'' while downstream implementations (C++ reference code, HLS kernels, vendor synthesis, and implementation) are assumed to be close enough. In practice, ``close enough'' is a moving target. A single change in numeric rounding, overflow behavior, or the ordering of edge updates can yield a program that still ``runs,'' still produces plausible outputs, and still passes superficial tests---while being observably different at the bit level.

This is not just a tooling inconvenience; it is a semantic gap. Hardware is unforgiving about hidden degrees of freedom: fixed-point arithmetic, saturation, scheduling constraints, and toolchain transformations decide the behavior that reaches silicon. Without an explicit semantic contract, there is no stable notion of correctness that can be checked across implementations, and there is no principled way to build an artifact that reviewers can independently validate. The result is a reproducibility failure mode that is particularly acute for sparse neural graphs and connectome-like topologies, where deterministic structure encourages hardware mapping but also amplifies the impact of ``minor'' schedule or numeric differences.

NEMA's design starts from the contract, not from the optimizer. The evaluated model is tick-based and deterministic: at each tick, the system applies a well-defined snapshot rule to compute updates on a sparse neural graph. The numeric model is fixed-point with explicit rounding (RNE) and saturating overflow, so that ``what happens'' is not an implementation detail. To make this contract testable, NEMA defines bit-exact equivalence as equality of per-tick SHA-256 digests over packed fixed-point state (e.g., Q8.8 voltages), turning correctness into a simple, automatable predicate.

We implement this contract as a small toolchain where the same NEMA program can be executed end-to-end through multiple representations. A manifest-driven harness runs a golden model, a C++ reference, and a hardware-oriented path (Vitis HLS csim/csynth, and---when applicable---Vivado implementation), producing machine-readable \texttt{bench\_report.json} outputs and consolidated paper tables. A separate \texttt{independent\_check} script re-runs the canonical checks for Paper A and fails the build if the expected pass/fail outcomes or key fields drift. The intent is deliberately narrow: every claim we make is traceable to a concrete artifact and can be verified by running a command.

What we do NOT claim (yet): this paper does not report on-board measurements of power, energy, or end-to-end latency on a physical FPGA; it does not claim an event-driven spiking model, LIF spiking dynamics, or biological plausibility results; and it does not claim Vivado implementation success for benchmarks where the artifact tables explicitly mark \texttt{vivado\_impl\_ok=false}. Throughput numbers in this draft are estimates derived from tool reports (e.g., II and timing) rather than on-board measurements.

Within that scope, the evaluation establishes three checkable facts. First, for the core set (B1/B2/B3), the verification harness reports mismatch=0 and \texttt{digestMatchOk=true}, and the independent checker reports success; these facts are recorded directly in the paper's bit-exactness table artifact. Second, for the same core set, the hardware pipeline completes through Vivado implementation and emits QoR/timing with non-null WNS; this is recorded directly in the QoR table artifact and corroborated by the Vivado coverage report. Third, the artifact pack itself is reproducible: the preflight command sequence (tests, independent check, and paper build) completes with exit code 0 and produces stable, hashed outputs; this is recorded by preflight logs and checksums.

Finally, NEMA is motivated by neuromorphic and connectome-inspired computation, but today's contribution is a deterministic, testable substrate on which such directions can be evaluated without semantic ambiguity. The external connectome-bundle workflow is exercised through the same manifest interface, demonstrating that the toolchain's notion of correctness and reproducibility extends beyond purely synthetic graphs, even when full hardware implementation success is not (yet) claimed for that benchmark.

What becomes possible once we have an FPGA board (measurement plan; no promised outcomes):
\begin{itemize}
\item Validate the boardless throughput estimates against measured clocking and observed ticks/sec on-device, using the same digest-based correctness predicate for outputs.
\item Measure power and energy per tick using board sensors or an external power meter, and record results in a new, versioned evidence artifact (power report + measurement protocol).
\item Demonstrate bitstream export and deploy, then run HW-in-the-loop tests that compare on-device digests to the golden reference for the core benchmarks.
\item Characterize variability across FPGA parts and builds (e.g., timing closure, achieved frequency) and report how it impacts the reproducible ``GO/NO-GO'' gates.
\end{itemize}

\subsection{Contributions}
\begin{itemize}
\item \textbf{Deterministic semantic contract with a bit-exact predicate.} NEMA specifies tick semantics plus a fixed-point contract and defines equivalence via per-tick SHA-256 digests; the bit-exact table shows mismatch=0 and \texttt{digestMatchOk=true} for B1/B2/B3 under this predicate (\texttt{papers/paperA/artifacts/tables/results\_bitexact.csv}).
\item \textbf{End-to-end, boardless hardware evidence with captured QoR/timing.} For B1/B2/B3, the toolchain records HLS csim/csynth success and Vivado implementation success with non-null WNS (\texttt{papers/paperA/artifacts/tables/results\_qor.csv}; \texttt{papers/paperA/artifacts/evidence/vivado\_coverage\_report.md}).
\item \textbf{Artifact-first reproducibility with independent verification.} The preflight suite (tests, independent claim check, and paper build) completes with exit 0 and logs are included for reviewers (\texttt{build/paperA\_preflight\_out\_v4/pytest\_q.*}; \texttt{build/paperA\_preflight\_out\_v4/independent\_check\_paperA.*}; \texttt{build/paperA\_preflight\_out\_v4/make\_clean\_paper.*}; \texttt{papers/paperA/artifacts/evidence/checksums.sha256}).
\item \textbf{Manifest-driven external graph workflow exercised.} The connectome-bundle path (B4) is executed through the same harness and produces \texttt{digestMatchOk=true} with a trace present, without claiming Vivado implementation success (\texttt{papers/paperA/artifacts/tables/results\_bitexact.csv} for B4; \texttt{build/bench\_verify\_b4/B4\_celegans\_external\_bundle/bench\_report.json}).
\end{itemize}

\paragraph{Claim-to-evidence map (explicit).}
\begin{enumerate}
\item[\textbf{C1}] \textbf{Bit-exactness on core benches (B1/B2/B3):} mismatch=0 and \texttt{digestMatchOk=true} over 20 ticks, independently re-checked.
\emph{Evidence:} \texttt{papers/paperA/artifacts/tables/results\_bitexact.csv} (rows B1/B2/B3: \texttt{mismatches\_len=0}, \texttt{digestMatchOk=true}, \texttt{independent\_check\_ok=true}); \texttt{build/paperA\_preflight\_out\_v4/independent\_check\_paperA.*}.
\item[\textbf{C2}] \textbf{Boardless hardware pipeline succeeds for core benches (B1/B2/B3):} HLS csim/csynth ok and Vivado implementation ok with non-null WNS captured.
\emph{Evidence:} \texttt{papers/paperA/artifacts/tables/results\_qor.csv} (rows B1/B2/B3: \texttt{hls\_csim\_ok=true}, \texttt{hls\_csynth\_ok=true}, \texttt{vivado\_impl\_ok=true}, numeric \texttt{wns}); \texttt{papers/paperA/artifacts/evidence/vivado\_coverage\_report.md}.
\item[\textbf{C3}] \textbf{Reproducible artifact pack and build:} tests, independent check, and paper build complete successfully; artifacts are hashed.
\emph{Evidence:} \texttt{build/paperA\_preflight\_out\_v4/pytest\_q.*}; \texttt{build/paperA\_preflight\_out\_v4/make\_clean\_paper.*}; \texttt{papers/paperA/artifacts/evidence/checksums.sha256}; \texttt{papers/paperA/artifacts/tables/results\_bitexact.csv}, \texttt{results\_qor.csv}, \texttt{results\_throughput.csv} (sha256 and timestamps listed in \texttt{review\_pack\_v4.md}).
\item[\textbf{C4}] \textbf{External connectome-bundle workflow is exercised through the same harness (without over-claiming hardware impl):} B4 verifies with mismatch=0 and \texttt{digestMatchOk=true}, while Vivado impl is explicitly not claimed as successful.
\emph{Evidence:} \texttt{papers/paperA/artifacts/tables/results\_bitexact.csv} (row B4: \texttt{verify\_ok=true}, \texttt{mismatches\_len=0}, \texttt{digestMatchOk=true}); \texttt{papers/paperA/artifacts/tables/results\_qor.csv} (row B4: \texttt{vivado\_impl\_ok=false}).
\end{enumerate}
